%
Et barn sitter i en huske. Husken trekkes bakover og slippes fra ro når setet er $h = \SI{1.2}{\metre}$ over det laveste punktet. Vi ser bort fra luftmotstand og friksjon.
\begin{enumerate}[a)]
    \item Hvor stor fart har barnet i det laveste punktet?
    \item Hvor høyt over det laveste punktet er barnet når farten er \SI{3.0}{\metre\per\second}?
    \item Snora i husken drar hele tiden i barnet, men gjør ikke noe arbeid. Hvorfor ikke? (Hint: hvilken vinkel er det mellom snorkraften og bevegelsesretningen?)
\end{enumerate}
